\begin{frame}{Meta-strategy}
\framesubtitle{군대에서 대회하기}

\begin{itemize}
  \item 가장 큰 어려움: 시간이 매우 부족합니다.
  \begin{itemize}
    \item 파견 전이었기 때문에 물리적으로 키보드를 잡을 시간이 절대적으로 없습니다.
    \item 불행 중 다행히 남는 일과시간을 논문을 읽고 설계하는데 쓸 수 있었습니다.
    \item 사지방에서 논문 프린트 \textrightarrow{} 남는 일과시간에 읽기, 설계 \textrightarrow{} 사지방에서 구현
  \end{itemize}
  \pause
  \item LLM은 확실히 Rust를 짤 때 오류가 적습니다.
  \begin{itemize}
    \item 순수 재미를 위해 Rust를 구현 언어로 선택했지만, 의외로 도움이 되었습니다.
    \item 의도가 코드에 명확하게 나타나고, 성능이 보장되면서 컴파일러가 오류를 잡아줍니다.
    \item LLM의 약점인 논리적 사고를 컴파일러가 대신 해 줄 수 있습니다.
  \end{itemize}
  \pause
  \item 아 이거 좋을것같은데 \textrightarrow{} 누군가 이미 했습니다.
  \begin{itemize}
    \item 특히 operator fusion 같은 유명한 주제는 많은 연구가 진행되어 있습니다.
    \item 선행 연구를 찾다가 아이디어를 얻을 때가 많이 있었습니다.
    \item 하지만 논문도 오류가 있을 수 있기 때문에 구현할 때 유의해야 합니다.
  \end{itemize}
\end{itemize}
\end{frame}